% ---
% id: receiver/rem:noworsefl
% kind: rem
% title: "The exact-window comparison under floating-point evaluation"
% summary: >-
%   Two steps of [[prop:exactnoworse]] compare evaluated numbers and survive an
%   evaluation; branch 1's use of the collapse lemma does not. There the
%   exact-window weight can sit below the analytic one by evaluation errors of
%   the bound function between the first block above the cutoff and one past
%   the depth.
% prominence: standard
% status: accepted
% uses: [prop:exactnoworse, def:pmc, lem:collapse]
% mentions: [thm:exact-and-evaluated, thm:emitted-bound]
% assumes: []
% refs: []
% provenance: newly-derived
% lean:
% ---

\begin{remark}
\label{rem:noworsefl}
Proposition~\ref{prop:exactnoworse} compares two \emph{expressions} in exact
arithmetic. Two of its three steps survive an evaluation term for term,
because they compare evaluated numbers with each other: $\ell(n)\ge\Bnd(n)$
is an evaluated maximum, and branch~2 minimises over exactly the evaluated
$\Bnd(n)$ that \eqref{eq:pmcdef} minimises over. The step that does not is
branch~1's use of Lemma~\ref{lem:collapse},
$\Bnd(n)\ge\Bnd(N+1)$ for $N+1<n\le m^{\ast}_t+1$: a statement about exact
values, and those further $\Bnd(n)$ are quantities \eqref{eq:pmcdef} never
evaluates, whose evaluated values can fall below the evaluated $\Bnd(N+1)$.
So on branch~1 the emitted value can sit below the one \eqref{eq:pmcdef}
would produce by the evaluation errors of those further values together with
that of $\Bnd(N+1)$, and this contract claims no bound on either
(Theorem~\ref{thm:exact-and-evaluated}\,(ii)). Unaffected either way are the
result that underwrites the emitted number (Theorem~\ref{thm:emitted-bound}) and
the fact that a declaration carrying no depth performs none of this
arithmetic.
\end{remark}
